\documentclass[12pt,a4paper]{ctexart}
\usepackage{amsmath,amssymb}
\usepackage{booktabs}
\usepackage{tabularx}
\usepackage{geometry}
\geometry{margin=2.5cm}
\usepackage{enumitem}
\usepackage{setspace}
\onehalfspacing
\usepackage{hyperref}
\usepackage{xcolor}

\pagestyle{plain} % 页码页末居中

\title{\bfseries 多重共线性的补救措施——极详细推导与讲解}
\author{}
\date{}

\begin{document}

\maketitle

\section*{引言}

在前面的学习中，我们已经知道了什么是多重共线性（解释变量之间高度相关），也学会了如何诊断它（方差扩大因子VIF、相关系数矩阵等）。诊断出病，接下来就要治。本文以初学者能跟上的细致程度，逐一讲解降低多重共线性的方法，重点推导岭回归法的数学原理。每一步代数变形、每一个直觉解释、每一个“为什么要这样做”都不省略。

\section*{第一步：我们要解决什么问题？}

想象一下，你要用一个人的“身高”和“体重”来预测他的“肺活量”。这两个变量（身高和体重）本身就高度相关——个子高的人通常也重。当它们同时出现在回归模型里时，回归系数就会变得极不稳定：今天换几个样本，身高的系数从正的变成负的，体重的系数也飘忽不定。这就是多重共线性带来的麻烦——\textbf{不是模型完全不能用，而是系数估计不可靠，标准误膨胀，t检验通不过}。

我们的目标是：用各种方法降低这种共线性，使得回归系数的估计变得更稳定、更可信。

\section{一、修正多重共线性的经验方法}

\subsection{方法1：剔除变量法}

\textbf{核心思路：} 既然共线性是因为某些变量“长得太像”引起的，那最直接的办法就是把多余的变量踢出去。

\textbf{具体操作：} 当回归方程中存在严重的多重共线性时，可以删除那个引起共线性的、且不太重要的变量。选择删谁时，要综合考虑三个信息：
\begin{itemize}
\item 回归系数的显著性检验（t检验的p值大不大？）
\item 方差扩大因子VIF（谁的VIF最大，谁就是“罪魁祸首”）
\item 经济含义（通过经济学常识判断哪个变量相对不重要）
\end{itemize}

\textcolor{red}{\textbf{⚠ 关键警告：}} 如果你不小心删掉了一个\textbf{重要变量}，就会产生\textbf{设定误差}（omitted variable bias）。这就像治病时把好细胞也一起杀掉了——共线性虽然消除了，但模型本身变错了，后果可能更严重。

\textbf{直觉比喻：} 你去医院查体，发现两个指标高度相关（比如血压和心率）。如果医生只保留一个指标来预测心脏病风险，这可以简化问题；但如果被删掉的那个恰好是更关键的指标，诊断就会出错。

\subsection{方法2：增大样本容量}

\textbf{核心思路：} 多重共线性本质上是\textbf{样本特性}，不是总体特性。换一批数据，共线性可能就没那么严重了。所以，尽可能多收集数据。

\textbf{为什么有效？数学推导如下：}

在前面学习方差扩大因子时，我们知道回归系数估计量 $\hat{\beta}_j$ 的方差为：
\[
\text{Var}(\hat{\beta}_j) = \frac{\sigma^2}{\sum x_i^2} \cdot \text{VIF}_j
\]
其中 $\sum x_i^2 = \sum(X_i - \bar{X})^2$ 是解释变量的离差平方和。

现在关键来了：当样本容量 $n$ 增加时，$\sum x_i^2$ 会怎样变化？
\begin{itemize}
\item 新增的观测值会带来新的 $(X_i - \bar{X})^2$ 项
\item 即使新的观测值恰好等于新平均值，$\sum x_i^2$ 也不会减小（最多保持不变）
\item 在大多数实际情况下，新数据会使得 $\sum x_i^2$ 增大
\end{itemize}

$\sum x_i^2$ 增大 → 分母增大 → $\text{Var}(\hat{\beta}_j)$ 减小 → 标准误减小 → 估计更精确

\textbf{直觉比喻：} 你用3个人的数据来估计“身高对体重的影响”，碰巧这3个人又高又瘦或又矮又胖，数据很混乱。但如果你收集了300人的数据，里面有各种各样的组合，共线性自然就弱了。

\textbf{局限性：} 有些数据（如宏观经济年度数据）很难增加样本量，一年才一个观测值，几十年也就几十个数据点。

\subsection{方法3：变换模型形式（差分法）}

\textbf{核心思路：} 原来的水平值之间高度相关，但如果对数据做一阶差分（本期减上期），差分后的变量之间的相关性通常会大大降低。

\textbf{数学推导——从水平模型到差分模型：}

原始模型（时间序列形式）：
\begin{equation}
Y_t = \beta_0 + \beta_1 X_{1t} + \beta_2 X_{2t} + \cdots + \beta_k X_{kt} + u_t \tag{4.22}
\end{equation}
其中下标 $t$ 表示第 $t$ 期。

写出第 $t-1$ 期的方程：
\[
Y_{t-1} = \beta_0 + \beta_1 X_{1,t-1} + \beta_2 X_{2,t-1} + \cdots + \beta_k X_{k,t-1} + u_{t-1}
\]

两式相减（本期减上期）：
\[
Y_t - Y_{t-1} = \beta_0 - \beta_0 + \beta_1(X_{1t} - X_{1,t-1}) + \beta_2(X_{2t} - X_{2,t-1}) + \cdots + \beta_k(X_{kt} - X_{k,t-1}) + (u_t - u_{t-1})
\]

注意：$\beta_0 - \beta_0 = 0$，截距项消失了！

定义差分符号 $\Delta$（读作“delta”），即 $\Delta Y_t = Y_t - Y_{t-1}$，$\Delta X_{jt} = X_{jt} - X_{j,t-1}$，$\Delta u_t = u_t - u_{t-1}$，得到：
\begin{equation}
\Delta Y_t = \beta_1 \Delta X_{1t} + \beta_2 \Delta X_{2t} + \cdots + \beta_k \Delta X_{kt} + \Delta u_t \tag{4.23}
\end{equation}

\textbf{为什么差分后共线性会减弱？}

举个具体例子：$X_1$（GDP）和 $X_2$（消费）的水平值高度相关，因为它们都是随时间持续增长的，趋势一致。但 $\Delta X_1$（GDP的增量）和 $\Delta X_2$（消费的增量）的相关性就弱得多——有的年份GDP涨得多但消费涨得少，有的年份反过来。

\textbf{EViews操作：} 在方程定义栏中输入 \texttt{Y-Y(-1) X1-X1(-1) X2-X2(-1) ... Xk-Xk(-1)} 即可得到 $\beta_1, \beta_2, \ldots, \beta_k$ 的估计值。

\textcolor{red}{\textbf{⚠ 注意事项：}}
\begin{itemize}
\item 差分会丢失信息（减少一个观测值，即样本从 $n$ 变成 $n-1$）
\item 差分模型的误差项 $\Delta u_t = u_t - u_{t-1}$ 可能存在\textbf{序列相关}，违背了经典假设
\item 因此，差分法虽然降低了共线性，但可能引入新的问题，使用时要慎重
\end{itemize}

\subsection{方法4：利用非样本先验信息}

\textbf{核心思路：} 如果经济学理论告诉我们某个参数是另一个参数的某个倍数，我们就可以利用这个约束来减少需要估计的参数个数，从而缓解共线性。

\textbf{数学推导：}

原模型：
\begin{equation}
Y_t = \beta_1 + \beta_2 X_{2t} + \beta_3 X_{3t} + u_t \tag{4.24}
\end{equation}

假设依据长期经验分析，我们知道 $\beta_3 = 0.2\beta_2$。

\textbf{这一步怎么用的？} 我们把 $\beta_3 = 0.2\beta_2$ 代入原模型：
\[
Y_t = \beta_1 + \beta_2 X_{2t} + 0.2\beta_2 X_{3t} + u_t
\]

把含 $\beta_2$ 的项合并：
\[
Y_t = \beta_1 + \beta_2(X_{2t} + 0.2X_{3t}) + u_t
\]

令 $X_t = X_{2t} + 0.2X_{3t}$（这是一个已知的新变量，因为我们知道 $X_{2t}$ 和 $X_{3t}$ 的数据，也知道系数0.2），则：
\begin{equation}
Y_t = \beta_1 + \beta_2 X_t + u_t \tag{4.25}
\end{equation}

\textbf{发生了什么？} 原来有3个参数要估计（$\beta_1, \beta_2, \beta_3$），现在只剩2个了（$\beta_1, \beta_2$），而且解释变量从2个合并成了1个！共线性问题直接消失了。

估计出 $\hat{\beta}_2$ 之后，$\hat{\beta}_3 = 0.2\hat{\beta}_2$ 也就自然得到了。

\textbf{直觉比喻：} 本来你要同时猜两个密码，它们互相关联导致你猜不准。现在有人告诉你第二个密码是第一个密码的0.2倍，你只需要猜一个就行了，问题大大简化。

\subsection{方法5：横截面数据与时序数据并用}

\textbf{核心思路：} 这是方法4（先验信息法）的一个变种。当我们无法直接获得先验信息时，可以用另一种数据来“替我们”估计部分参数。

\textbf{完整推导——以中国家用轿车需求为例：}

设定模型：
\begin{equation}
\ln Y_t = \beta_1 + \beta_2 \ln P_t + \beta_3 \ln I_t + u_t \tag{4.26}
\end{equation}
\begin{itemize}
\item $Y_t$：轿车销售量
\item $P_t$：平均价格
\item $I_t$：消费者收入
\item $\beta_2$：价格弹性（价格变化1\%，需求变化百分之几）
\item $\beta_3$：收入弹性（收入变化1\%，需求变化百分之几）
\end{itemize}

\textbf{问题：} 在时间序列数据中，价格 $P_t$ 和收入 $I_t$ 通常高度相关（两者都随时间持续增长），导致 $\beta_2$ 和 $\beta_3$ 估计不准。

\textbf{托宾的解决办法，分三步走：}

\textbf{第①步：} 用横截面数据（比如城镇/农村住户调查数据）估计收入弹性 $\beta_3$。横截面数据是在同一时间点上、不同家庭的数据，此时各家庭的收入差异大但价格基本相同，所以不会受价格和收入共线性的困扰。设估计结果为 $\hat{\beta}_3^*$。

\textbf{第②步：} 将 $\hat{\beta}_3^*$ 代入原时间序列模型。将式(4.26)改写：
\[
\ln Y_t - \hat{\beta}_3^* \ln I_t = \beta_1 + \beta_2 \ln P_t + u_t
\]
令 $Y_t^* = \ln Y_t - \hat{\beta}_3^* \ln I_t$，则：
\begin{equation}
Y_t^* = \beta_1 + \beta_2 \ln P_t + u_t \tag{4.27}
\end{equation}

\textbf{第③步：} 现在只需要用时间序列数据估计一元回归(4.27)，得到 $\hat{\beta}_2$（价格弹性）。

\textbf{发生了什么？} 原来二元回归中的共线性问题，被拆成了两步解决：横截面估计一个参数，时间序列估计另一个参数。

\textcolor{red}{\textbf{⚠ 关键假设：}} 收入弹性的横截面估计 $\hat{\beta}_3^*$ 和从纯粹时间序列分析中得到的估计是一样的。如果这个假设不成立（比如横截面上的收入效应和随时间变化的收入效应不同），结果就会有偏差。

\subsection{方法6：变量变换}

有时不改变模型形式，只对变量本身做变换，也能降低共线性。常用的变换方式有四种：

\textbf{（1）计算相对指标}

原来是总量指标（如GDP总量、消费总量），改为人均指标（人均GDP、人均消费）或结构相对数（消费占GDP的比重）。经过这样处理，原来同趋势增长的总量变成了相对稳定的比率，共线性可能降低。

\textbf{（2）将名义数据转换为实际数据}

名义数据包含价格变动和物量变动双重影响。剔除价格影响后的实际数据只反映物量变化，更可能揭示变量之间真实的数量关系，从而降低共线性。

\textbf{（3）将小类指标合并成大类指标}

例如，工业增加值和建筑业增加值通常高度相关（因为它们都受宏观经济景气影响），同时引入模型会产生严重共线性。将它们合并为“第二产业增加值”，用一个变量代替两个，共线性自然消除。

\textbf{（4）将总量指标进行对数变换}

原来的线性模型 $Y = \beta_0 + \beta_1 X_1 + \beta_2 X_2 + u$，对变量取对数后变成双对数模型 $\ln Y = \beta_0 + \beta_1 \ln X_1 + \beta_2 \ln X_2 + u$，此时系数的含义从“X变化1个单位，Y变化多少”变成“X变化1\%，Y变化百分之几”。增长率之间的关系往往比水平值之间的关系更独立，共线性可能降低。

\textcolor{red}{\textbf{⚠ 注意：}} 变量变换只是一种尝试性方法，有时效果很好，有时可能没有效果。不能保证每次都成功。

\section{二、逐步回归法}

\subsection{核心思路}

逐步回归法的想法很朴素：既然多个变量一起放进去会共线性，那我就一个一个往里加，每加一个都检查它是否“捣乱”。

\subsection{具体步骤（逐步拆解）}

\textbf{步骤1：} 用被解释变量 $Y$ 对每一个候选解释变量 $X_1, X_2, \ldots, X_k$ 分别做简单回归（一元回归），一共做 $k$ 个回归。

\textbf{步骤2：} 比较这 $k$ 个回归的结果，选出对 $Y$ 贡献最大的那个解释变量（看哪个回归的 $\bar{R}^2$ 最大）。把这个变量作为“基础变量”，它所对应的回归方程作为“基础方程”。

\textbf{步骤3：} 在基础方程的基础上，逐个引入其余的解释变量。每次引入一个新变量后，检查以下三个指标的变化：

\begin{table}[htbp]
\centering
\caption{逐步回归中引入新变量后的判断准则}
\begin{tabularx}{\textwidth}{@{}l *{4}{>{\centering\arraybackslash}X}@{}}
\toprule
\textbf{情形} & \textbf{$\bar{R}^2$} & \textbf{$F$检验} & \textbf{其他原有变量的$t$检验} & \textbf{判断} \\
\midrule
① & 改善了 & 通过 & 仍然显著 & $\checkmark$ 保留该变量 \\
② & 没改善 & 没改善 & 没受影响 & $\times$ 该变量多余 \\
③ & 没改善 & 没改善 & 显著受影响（系数变化大或不再显著） & $\times$ 出现严重共线性 \\
\bottomrule
\end{tabularx}
\end{table}

\textbf{对情形③的详细解释：} 为什么新变量进来后，原来显著的变量突然不显著了？因为新变量和原来的某个变量高度相关（共线性！），它们“争抢”解释力，导致系数估计变得不稳定。这时应该把新变量剔除。

\textbf{步骤4：} 重复步骤3，直到所有候选变量都尝试过。最终保留在模型中的变量，既对 $Y$ 有重要解释贡献，彼此之间又没有严重共线性。

\subsection{优点与局限}

\textbf{优点：} 最后保留的变量是统计上显著的、彼此间共线性不明显的，模型简洁有效。

\textcolor{red}{\textbf{⚠ 局限：}} 逐步回归可能删除了重要的相关变量，导致\textbf{设定偏误}。这就像医生在排除干扰因素时，不小心把真正的病因也排除了，治标不治本。

\section{三、岭回归法}

这是本节最重要的数学内容。前面几种方法都是“经验方法”，逻辑上比较直观；岭回归法则需要对数学公式进行严谨推导。

\subsection{1. 岭回归的含义——为什么需要它？}

\textbf{问题回顾：} 当解释变量之间存在多重共线性时，矩阵 $X'X$ 接近奇异（即行列式 $|X'X| \approx 0$）。我们知道OLS估计量为：
\[
\hat{\beta} = (X'X)^{-1}X'Y
\]
而 $\hat{\beta}$ 的方差-协方差矩阵为：
\[
\text{Var}(\hat{\beta}) = \sigma^2(X'X)^{-1}
\]

\textbf{关键推理链条（无跳步）：}

① $|X'X| \approx 0$（因为多重共线性使 $X'X$ 接近奇异）

② $(X'X)^{-1}$ 中的元素会变得非常大（因为逆矩阵的元素涉及 $1/|X'X|$ 的运算，分母很小，值就很大）

③ $\text{Var}(\hat{\beta}) = \sigma^2(X'X)^{-1}$ 中的元素因此也非常大

④ 方差很大 → 标准误很大 → t统计量很小 → 系数不显著

\textbf{这就是多重共线性导致OLS估计“不稳定”的数学根源！}

\textbf{岭回归的核心思想：} 既然问题出在 $X'X$ 接近奇异，那我就人为地给 $X'X$ 加上一个正常数对角矩阵 $kI$（$k > 0$），使得 $|X'X + kI|$ 远离0。

\textbf{直觉比喻：} 想象一个薄薄的冰面（$X'X$），你站在上面随时可能裂开（行列式接近0）。现在在冰面上铺一层木板（$kI$），冰面变厚了（$X'X + kI$），承重能力增强了（行列式远离0），你就安全了。

\subsubsection{岭回归估计量的定义}
\begin{equation}
\hat{\beta}(k) = (X'X + kI)^{-1}X'Y \tag{4.28}
\end{equation}
其中：
\begin{itemize}
\item $\hat{\beta}(k)$：称为 $\beta$ 的\textbf{岭回归估计量}（ridge estimate）
\item $k$：\textbf{岭回归参数}，也叫岭参数或偏倚参数，$k > 0$
\item $I$：单位矩阵（对角线上全是1，其余全是0的方阵）
\item $kI$：对角矩阵，对角线上全是 $k$，其余全是0
\end{itemize}

\textbf{极端情况的理解：}
\begin{itemize}
\item 当 $k = 0$ 时：$\hat{\beta}(0) = (X'X)^{-1}X'Y = \hat{\beta}$，岭回归退化为普通OLS估计
\item 当 $k$ 很小时：$\hat{\beta}(k)$ 接近OLS估计，但可能仍然不稳定
\item 当 $k$ 逐渐增大时：$X'X + kI$ 越来越远离奇异，$\hat{\beta}(k)$ 趋于稳定
\item 当 $k \to \infty$ 时：$(X'X + kI)^{-1} \approx (kI)^{-1} = \frac{1}{k}I$，$\hat{\beta}(k) \to 0$，所有系数都被压缩到0
\end{itemize}

所以，\textbf{选择合适的 $k$ 值是岭回归的关键}：$k$ 太小没效果，$k$ 太大系数全被压没了。

\subsection{2. 岭回归估计的性质（极其详细的推导）}

\subsubsection{性质1：岭回归的参数估计是回归参数的有偏估计}

\textbf{目标：} 证明 $E[\hat{\beta}(k)] \neq \beta$（当 $k \neq 0$ 时）

\textbf{推导（无跳步）：}

\[
E[\hat{\beta}(k)] = E[(X'X + kI)^{-1}X'Y]
\]

这里，$(X'X + kI)^{-1}$ 和 $X'$ 都是常数矩阵（因为 $X$ 被假定为固定），只有 $Y$ 是随机的。根据期望的线性性质，常数矩阵可以提到期望符号外面：
\[
= (X'X + kI)^{-1}X' \cdot E[Y]
\]

现在关键是 $E[Y]$ 等于什么。由真实模型 $Y = X\beta + u$，且 $E(u) = 0$，所以：
\[
E[Y] = E[X\beta + u] = X\beta + E(u) = X\beta + 0 = X\beta
\]

代入回去：
\[
= (X'X + kI)^{-1}X' \cdot X\beta = (X'X + kI)^{-1}(X'X)\beta \tag{4.29}
\]

\textbf{分析这个结果：}
\begin{itemize}
\item 当 $k = 0$ 时：$(X'X + 0 \cdot I)^{-1}(X'X) = (X'X)^{-1}(X'X) = I$，所以 $E[\hat{\beta}(0)] = I\beta = \beta$。$\checkmark$ 无偏！
\item 当 $k \neq 0$ 时：$(X'X + kI)^{-1} \neq (X'X)^{-1}$，所以 $(X'X + kI)^{-1}(X'X) \neq I$，从而 $E[\hat{\beta}(k)] \neq \beta$。$\times$ 有偏！
\end{itemize}

\textbf{有偏性的直觉理解：} 岭回归为了获得更稳定的估计（减小方差），付出了代价——引入了偏差。这就像射击时，你宁可枪口稍微偏一点但稳定性很高（每次都打在同一个地方），也不愿枪口对准靶心但弹着点到处散开。

\subsubsection{性质2：$\hat{\beta}(k)$ 是最小二乘估计的一个线性变换}

\textbf{推导：}

已知OLS估计量 $\hat{\beta} = (X'X)^{-1}X'Y$。

岭估计量 $\hat{\beta}(k) = (X'X + kI)^{-1}X'Y$。

我们想把 $\hat{\beta}(k)$ 表示成 $\hat{\beta}$ 的线性函数：
\[
\hat{\beta}(k) = (X'X + kI)^{-1}X'Y
\]

注意 $X'Y = X'X \cdot (X'X)^{-1}X'Y = X'X \cdot \hat{\beta}$（因为 $\hat{\beta} = (X'X)^{-1}X'Y$，两边左乘 $X'X$ 即得）。

代入：
\[
\hat{\beta}(k) = (X'X + kI)^{-1}(X'X)\hat{\beta}
\]

记 $W = (X'X + kI)^{-1}(X'X)$，则 $\hat{\beta}(k) = W\hat{\beta}$。

因为 $W$ 是常数矩阵（仅依赖于 $X$ 和 $k$），所以 $\hat{\beta}(k)$ 是 $\hat{\beta}$ 的\textbf{线性变换}。

同时，由于 $\hat{\beta}$ 本身是 $Y$ 的线性函数，$\hat{\beta}(k)$ 也是 $Y$ 的线性函数。

\subsubsection{性质3：岭回归估计量的方差比OLS估计量的方差小}

\textbf{目标：} 证明 $\text{Var}(\hat{\beta}(k))$ 比 $\text{Var}(\hat{\beta})$ 小

\textbf{推导（无跳步）：}

由性质2，$\hat{\beta}(k) - E[\hat{\beta}(k)] = W(\hat{\beta} - \beta)$（此处为简化说明，严格推导涉及偏倚项，但我们关注方差部分）。

更直接地，我们从定义出发：
\[
\hat{\beta}(k) = (X'X + kI)^{-1}X'Y = (X'X + kI)^{-1}X'(X\beta + u)
\]
\[
= (X'X + kI)^{-1}X'X\beta + (X'X + kI)^{-1}X'u
\]

第一项是常数（因为 $X$ 和 $\beta$ 固定），第二项含有随机项 $u$。所以 $\hat{\beta}(k)$ 的波动完全来自第二项。

\[
\text{Var}(\hat{\beta}(k)) = \text{Var}[(X'X + kI)^{-1}X'u]
\]

因为 $(X'X + kI)^{-1}X'$ 是常数矩阵，根据 $\text{Var}(Au) = A\text{Var}(u)A'$：
\[
= (X'X + kI)^{-1}X' \cdot \sigma^2 I \cdot X(X'X + kI)^{-1}
\]
\[
= \sigma^2(X'X + kI)^{-1}X'X(X'X + kI)^{-1}
\]

对比OLS的方差：$\text{Var}(\hat{\beta}) = \sigma^2(X'X)^{-1}$

\textbf{为什么岭回归方差更小？直觉解释：}

$(X'X + kI)$ 的特征值比 $X'X$ 的特征值大（每个特征值都加了 $k$）。矩阵逆的特征值是原特征值的倒数，所以 $(X'X + kI)^{-1}$ 的特征值比 $(X'X)^{-1}$ 的特征值小。方差-协方差矩阵的“总大小”（用迹或行列式衡量）因此减小了。

\textbf{简单比喻：} OLS方差像是站在一个很窄的柱子上（不稳定，晃动大）；岭回归给柱子加宽了底座（加了 $kI$），站上去更稳了（方差小了），但站的位置偏了（有偏了）。

\subsection{3. 岭回归参数 $k$ 的选择}

\textbf{核心矛盾：} $k$ 越大 → 偏倚越大、方差越小；$k$ 越小 → 偏倚越小、方差越大。

\textbf{用什么准则来权衡？} 最小均方误差（MSE）原则。

均方误差的定义：
\[
\text{MSE}(\hat{\beta}(k)) = E[\hat{\beta}(k) - \beta][\hat{\beta}(k) - \beta]' = \text{方差} + \text{偏倚}^2
\]
即：MSE = 方差部分 + 偏倚的平方部分。

我们要找使 $\text{MSE}(\hat{\beta}(k))$ 达到最小的 $k^*$。

\textbf{问题：} 最优 $k^*$ 依赖于未知参数 $\beta$ 和 $\sigma^2$，所以理论上最优的 $k$ 是无法直接计算的。实际中必须通过样本数据来确定。

\textbf{常用方法：}

\begin{table}[htbp]
\centering
\caption{岭回归参数 $k$ 的选择方法}
\begin{tabularx}{\textwidth}{@{}l X@{}}
\toprule
\textbf{方法} & \textbf{基本思路} \\
\midrule
岭迹法 & 画出 $\hat{\beta}_j(k)$ 随 $k$ 变化的曲线（岭迹图），选择系数开始趋于稳定时的 $k$ \\
方差扩大因子法 & 选择使所有VIF都降到合理范围（如小于10）的 $k$ \\
残差平方和法 & 在残差平方和不过分增大的前提下，选尽可能大的 $k$ \\
\bottomrule
\end{tabularx}
\end{table}

\textbf{逐步搜索法的实际操作：}
\begin{enumerate}
\item 从很小的 $k$ 开始（如 $k = 0.01$）
\item 逐渐增大 $k$（如 0.02, 0.05, 0.1, 0.2, 0.5, 1, 2, \dots）
\item 每次计算 $\hat{\beta}(k)$，观察系数值的变化
\item 当系数值开始趋于稳定时，选择该 $k$ 值
\end{enumerate}

\textcolor{red}{\textbf{⚠ 局限：}} 逐步搜索法确定 $k$ 值缺乏严密的理论依据，具有一定的主观性，是一种将定性分析与定量分析相结合的方法。

\subsection*{拓展：与岭回归相似的其他方法}

\begin{table}[htbp]
\centering
\caption{与岭回归相似的其他方法}
\begin{tabularx}{\textwidth}{@{}l X@{}}
\toprule
\textbf{方法} & \textbf{核心特点} \\
\midrule
Lasso回归 & 用 $\lambda\sum|\beta_j|$ 替代岭回归的 $\lambda\sum\beta_j^2$，可以将某些系数压缩到恰好0（自动变量选择） \\
适应性Lasso回归 & 对Lasso的惩罚权重做自适应调整，改善Lasso的某些理论缺陷 \\
主成分分析 & 将原来的共线性变量转换为少数几个互不相关的主成分，用主成分做回归 \\
偏最小二乘回归 & 类似主成分但兼顾与因变量的相关性，提取成分时同时考虑自变量和因变量 \\
\bottomrule
\end{tabularx}
\end{table}

\section*{总结：整篇讲解的逻辑链条}

\begin{enumerate}
\item \textbf{诊断出共线性后怎么办？} → 六种经验方法（剔除变量、增大样本、差分法、先验信息、数据并用、变量变换）+ 逐步回归法 + 岭回归法
\item \textbf{经验方法的共同思路} → 减少需要估计的参数个数，或降低变量之间的相关性
\item \textbf{逐步回归} → 一个个变量试探，保留有用的、剔除“捣乱”的
\item \textbf{岭回归} → 核心数学：给 $X'X$ 加 $kI$ 使其远离奇异，代价是估计量从无偏变为有偏，但方差大幅减小；用MSE准则权衡偏倚与方差
\item \textbf{岭参数 $k$ 的选择} → 没有公认最优方法，岭迹法、VIF法、逐步搜索法都是实用但带主观性的方法
\end{enumerate}

\end{document}